% Zdroj: PDF priklady-jaro-2023.pdf, fyzická str. 2. Značka: starší neznačená sada. Zařazeno: M2.
\section{Lineární zobrazení}
\szzotazka{jaro 2023}{3}{Lineární zobrazení}

\noindent\textbf{Zadání.}\par
\begin{enumerate}
  \item Definujte \textit{matici lineárního zobrazení vzhledem k zadaným bázím}.
  \item Formulujte tvrzení o matici složeného zobrazení, tj.\ jak se matice lineárního zobrazení mění při skládání zobrazení.
  \item Uvažme lineární zobrazení $g : \mathbb{Z}_5^2 \to \mathbb{Z}_5^2$ splňující $f \circ g \circ h = \mathrm{id}$, přičemž lineární zobrazení $f$ a $h$ jsou dána následovně:
    \[
      f((1,0)^T) = h((0,1)^T) = (3,4)^T,
    \]
    \[
      f((0,1)^T) = h((1,0)^T) = (4,1)^T.
    \]
    Najděte matici $g$ vzhledem ke kanonickým bázím.
\end{enumerate}

\medskip
\noindent\textbf{Náčrt řešení.}\par
\textit{(V originálním PDF bez náčrtu; doplněno vzorové řešení.)}
\begin{enumerate}
  \item \textit{Matice lineárního zobrazení.} Buď $f:U\to V$ lineární, $B_U=\{x_1,\dots,x_n\}$ báze $U$ a $B_V$ báze $V$. \emph{Maticí $f$ vzhledem k~bázím} $B_U,B_V$ (značíme ${}_{B_V}[f]_{B_U}$) rozumíme matici, jejíž $j$-tý sloupec tvoří souřadnice obrazu $[f(x_j)]_{B_V}$. Splňuje $[f(x)]_{B_V}={}_{B_V}[f]_{B_U}\cdot[x]_{B_U}$.
  \item \textit{Matice složeného zobrazení.} Pro $f:U\to V$, $g:V\to W$ a~báze $B_U,B_V,B_W$ platí
    \[
      {}_{B_W}[g\circ f]_{B_U}={}_{B_W}[g]_{B_V}\cdot{}_{B_V}[f]_{B_U}.
    \]
    Skládání zobrazení tedy odpovídá násobení jejich matic.
  \item Sloupce matice zobrazení vzhledem ke kanonickým bázím jsou obrazy kanonických vektorů, takže (počítáme v~tělese $\mathbb{Z}_5$)
    \[
      F=\begin{pmatrix}3&4\\4&1\end{pmatrix},\qquad H=\begin{pmatrix}4&3\\1&4\end{pmatrix},
    \]
    kde $F_{*1}=f((1,0)^T)=(3,4)^T$, $F_{*2}=f((0,1)^T)=(4,1)^T$ a~obdobně $H_{*1}=h((1,0)^T)=(4,1)^T$, $H_{*2}=h((0,1)^T)=(3,4)^T$. Z~rovnosti $f\circ g\circ h=\mathrm{id}$ plyne $FGH=I$, tedy $G=F^{-1}H^{-1}$.

    Inverzi matice řádu~2 počítáme jako $\begin{psmallmatrix}a&b\\c&d\end{psmallmatrix}^{-1}=(\det)^{-1}\begin{psmallmatrix}d&-b\\-c&a\end{psmallmatrix}$ (vše modulo~5):
    \begin{itemize}
      \item $\det F=3\cdot1-4\cdot4=-13\equiv2$, $\ 2^{-1}=3$, takže $F^{-1}=3\begin{psmallmatrix}1&-4\\-4&3\end{psmallmatrix}=\begin{psmallmatrix}3&3\\3&4\end{psmallmatrix}$;
      \item $\det H=4\cdot4-3\cdot1=13\equiv3$, $\ 3^{-1}=2$, takže $H^{-1}=2\begin{psmallmatrix}4&-3\\-1&4\end{psmallmatrix}=\begin{psmallmatrix}3&4\\3&3\end{psmallmatrix}$.
    \end{itemize}
    Pak
    \[
      G=F^{-1}H^{-1}=\begin{pmatrix}3&3\\3&4\end{pmatrix}\begin{pmatrix}3&4\\3&3\end{pmatrix}=\begin{pmatrix}3&1\\1&4\end{pmatrix}\pmod 5.
    \]
    Ověření: $FG=\begin{psmallmatrix}3&4\\3&3\end{psmallmatrix}$ a~$(FG)H=I$, takže $f\circ g\circ h=\mathrm{id}$.
\end{enumerate}
